% Standalone Appendix A note for the book.
% This file is intended to be pasted into the appendix text or included with
% % Standalone Appendix A note for the book.
% This file is intended to be pasted into the appendix text or included with
% % Standalone Appendix A note for the book.
% This file is intended to be pasted into the appendix text or included with
% % Standalone Appendix A note for the book.
% This file is intended to be pasted into the appendix text or included with
% \input{APPENDIX_A_PROCESS_GLOBAL_LIMITATIONS}.

\paragraph{Implementation note: per-node versus process-global state.}
The current SWI-Prolog proof-of-concept supports running multiple Web Prolog
nodes inside the same SWI process, and the request-scoped runtime state is now
largely isolated per node. Profile selection, authentication mode, sandbox
mode, node URL, principal policy, shared database, timeout policy, and the
stateless \texttt{/call} continuation cache are all resolved per node during
ordinary HTTP and WebSocket execution.

\paragraph{What still remains process-global?}
The remaining global state concerns only \emph{no-context local convenience
use} inside one SWI process. Two examples are especially relevant:

\begin{itemize}
\item the fallback shared-database facade used by direct local calls outside a
request context, and
\item the fallback self-node registration used when no per-node request context
is active.
\end{itemize}

In practice this means that if several nodes are started inside one SWI
process, then direct local interactions that do not run under a node-specific
request or actor context still follow the most recently started node.

\paragraph{Practical scope of the limitation.}
This limitation matters only for in-process local experimentation. The actual
node-facing HTTP and WebSocket APIs now keep their shared database, timeout,
and continuation-cache state separate per node, even when multiple nodes are
hosted in one SWI-Prolog process.
.

\paragraph{Implementation note: per-node versus process-global state.}
The current SWI-Prolog proof-of-concept supports running multiple Web Prolog
nodes inside the same SWI process, and the request-scoped runtime state is now
largely isolated per node. Profile selection, authentication mode, sandbox
mode, node URL, principal policy, shared database, timeout policy, and the
stateless \texttt{/call} continuation cache are all resolved per node during
ordinary HTTP and WebSocket execution.

\paragraph{What still remains process-global?}
The remaining global state concerns only \emph{no-context local convenience
use} inside one SWI process. Two examples are especially relevant:

\begin{itemize}
\item the fallback shared-database facade used by direct local calls outside a
request context, and
\item the fallback self-node registration used when no per-node request context
is active.
\end{itemize}

In practice this means that if several nodes are started inside one SWI
process, then direct local interactions that do not run under a node-specific
request or actor context still follow the most recently started node.

\paragraph{Practical scope of the limitation.}
This limitation matters only for in-process local experimentation. The actual
node-facing HTTP and WebSocket APIs now keep their shared database, timeout,
and continuation-cache state separate per node, even when multiple nodes are
hosted in one SWI-Prolog process.
.

\paragraph{Implementation note: per-node versus process-global state.}
The current SWI-Prolog proof-of-concept supports running multiple Web Prolog
nodes inside the same SWI process, and the request-scoped runtime state is now
largely isolated per node. Profile selection, authentication mode, sandbox
mode, node URL, principal policy, shared database, timeout policy, and the
stateless \texttt{/call} continuation cache are all resolved per node during
ordinary HTTP and WebSocket execution.

\paragraph{What still remains process-global?}
The remaining global state concerns only \emph{no-context local convenience
use} inside one SWI process. Two examples are especially relevant:

\begin{itemize}
\item the fallback shared-database facade used by direct local calls outside a
request context, and
\item the fallback self-node registration used when no per-node request context
is active.
\end{itemize}

In practice this means that if several nodes are started inside one SWI
process, then direct local interactions that do not run under a node-specific
request or actor context still follow the most recently started node.

\paragraph{Practical scope of the limitation.}
This limitation matters only for in-process local experimentation. The actual
node-facing HTTP and WebSocket APIs now keep their shared database, timeout,
and continuation-cache state separate per node, even when multiple nodes are
hosted in one SWI-Prolog process.
.

\paragraph{Implementation note: per-node versus process-global state.}
The current SWI-Prolog proof-of-concept supports running multiple Web Prolog
nodes inside the same SWI process, and the request-scoped runtime state is now
largely isolated per node. Profile selection, authentication mode, sandbox
mode, node URL, principal policy, shared database, timeout policy, and the
stateless \texttt{/call} continuation cache are all resolved per node during
ordinary HTTP and WebSocket execution.

\paragraph{What still remains process-global?}
The remaining global state concerns only \emph{no-context local convenience
use} inside one SWI process. Two examples are especially relevant:

\begin{itemize}
\item the fallback shared-database facade used by direct local calls outside a
request context, and
\item the fallback self-node registration used when no per-node request context
is active.
\end{itemize}

In practice this means that if several nodes are started inside one SWI
process, then direct local interactions that do not run under a node-specific
request or actor context still follow the most recently started node.

\paragraph{Practical scope of the limitation.}
This limitation matters only for in-process local experimentation. The actual
node-facing HTTP and WebSocket APIs now keep their shared database, timeout,
and continuation-cache state separate per node, even when multiple nodes are
hosted in one SWI-Prolog process.
